\section{Current Situation}

\subsection{Problem Statement}

The current situation is marked by a vicious continuous cycle: the framework created to monitor and prevent illicit financial flows has increased greatly over the last 30 years, however, the volume of laundered money has also increased.\footnote{United Nations Office on Drugs and Crime. (2011). \emph{Estimating illicit financial flows resulting from drug trafficking and other transnational organized crimes}. UNODC.} Based on UNODC estimates, 2-5\% of global GDP (anywhere from \$800 billion to \$2 trillion) is laundered annually through the international financial system.\footnote{United Nations Office on Drugs and Crime. (2011). \emph{Estimating illicit financial flows resulting from drug trafficking and other transnational organized crimes}. UNODC.} The fundamental problem is that as organisations continue to develop ways to combat crime and illicit funding mechanisms, the people using those funding mechanisms discover weaknesses in the funding mechanisms.\footnote{Levi, M. (2020). Making sense of professional enablers' involvement in laundering organized crime proceeds. \emph{Trends in Organized Crime, 24}(1), 96–110.}

Outflows of illicit funds disproportionately affect the states with the weakest regulatory infrastructure. This indicates that the countries which are most damaged by illicit finance are also the ones least able to combat it.\footnote{Ndikumana, L., \& Boyce, J. K. (2011). \emph{Africa's odious debts: How foreign loans and capital flight bled a continent}. Zed Books.}

Traditional methods of money laundering include using shell companies, exploiting relationships between correspondent banks and pushing money through real estate. All these methods remain prevalent.\footnote{Palan, R., Murphy, R., \& Chavagneux, C. (2010). \emph{Tax havens: How globalization really works}. Cornell University Press.} However, there have been significant developments in financial technologies which have allowed for more sophisticated channels. This includes cryptocurrencies and decentralised finance platforms (DeFi), which operate outside the traditional banking system and its regulations. This indicates that states and organisations struggle to monitor the use of these new technologies and struggle to monitor the flow of all these illicit funds.\footnote{Financial Action Task Force. (2021). \emph{Updated guidance for a risk-based approach to virtual assets and virtual asset service providers}. \url{https://www.fatf-gafi.org/}} The largest method of money laundering remains trade-based. This includes manipulating the price, quantity, or quality of goods during international trade transactions. It is also one of the most difficult to detect.\footnote{Zdanowicz, J. S. (2009). Trade-based money laundering and terrorist financing. \emph{Review of Law and Economics, 5}(2), 855–878.}

One of the main deficiencies of the money laundering regulations, is that it depends on political commitment of members rather than binding international laws. Look no further than the Financial Action Task Force.\footnote{Nance, M. T. (2018). The regime that FATF built: An introduction to the Financial Action Task Force. \emph{Crime, Law and Social Change, 69}(2), 109–129.} It has a mutual evaluation process, which has often been criticised for applying the standards unevenly. Developing states are often placed on grey or blacklists when they have regulatory gaps, whereas countries in the West do not comparable consequences.\footnote{Cobham, A., Janský, P., \& Meinzer, M. (2015). The Financial Secrecy Index: Shedding new light on the geography of secrecy. \emph{Economic Geography, 91}(3), 281–303.}

INTERPOL suffers from constraints on its ability to coordinate multi-jurisdictional financial investigations. It is constrained by the requirement for voluntary cooperation from national law enforcement agencies.\footnote{Deflem, M. (2010). \emph{The policing of terrorism: Organizational and global perspectives}. Routledge.} These national law enforcement agencies all operate under different legal frameworks, data protection regimes, and often have different political agendas.\footnote{Sheptycki, J. (2004). Organizational pathologies in police intelligence systems: Some contributions to the lexicon of intelligence-led policing. \emph{European Journal of Criminology, 1}(3), 307–332.}

\subsection{Past Actions}

The first phase towards establishing the Financial Action Task Force was by the G7 at the 1989 Paris Summit. They intended to create international standards to combat money laundering.\footnote{Financial Action Task Force. (1990). \emph{The forty recommendations}. FATF/OECD.} They created 40 recommendations with a basic framework which included\footnote{Gilmore, W. C. (2004). \emph{Dirty money: The evolution of international measures to counter money laundering and the financing of terrorism} (3rd ed.). Council of Europe Publishing.}:

\begin{itemize}
  \item Criminalising money laundering
  \item Customer due diligence
  \item Obligations for financial institutions
  \item Requirements for suspicious transaction reporting
\end{itemize}

Then in 1988 the United Nations Convention against Illicit Traffic in Narcotic Drugs and Psychotropic Substances, also known as the Vienna Convention, established an international legal basis for criminalising the laundering of money coming from drug operations.\footnote{United Nations. (1988). \emph{United Nations Convention against Illicit Traffic in Narcotic Drugs and Psychotropic Substances}. \url{https://www.unodc.org/}} The 2000 Convention against Transnational Organized Crime, also known as the Palermo Convention, was used to broaden the list of offences to include almost all serious criminal activity.\footnote{United Nations. (2000). \emph{United Nations Convention against Transnational Organized Crime}. \url{https://www.unodc.org/}} The 2003 United Nations Convention against Corruption further extended the framework to the laundering of money from bribery and embezzlement of public officials.\footnote{United Nations. (2003). \emph{United Nations Convention against Corruption}. \url{https://www.unodc.org/}}

After 9/11, the framework expanded to include the financing of terrorism. 8 special recommendations were created (called the Eight Special Recommendations on Terrorist Financing). And the United Nations Security Council adopted Resolution 1373 which required all member states to criminalise terrorist financing and to freeze terrorist-related assets.\footnote{United Nations Security Council. (2001). \emph{Resolution 1373} (S/RES/1373). \url{https://undocs.org/S/RES/1373(2001)}} \footnote{Winer, J. M. (2008). Countering terrorist finance: A work, mostly in progress. \emph{The ANNALS of the American Academy of Political and Social Science, 618}(1), 112–132.} This also strengthened the tension between surveillance of financials and the civil liberties potentially being encroached upon.\footnote{de Goede, M. (2012). \emph{Speculative security: The politics of pursuing terrorist monies}. University of Minnesota Press.}

INTERPOL has increased its role into financial crime enforcement, for example, the creation of the Financial Crimes and Anti-Corruption Centre (IFCACC). They coordinate multi-jursisdictional investigations and provide analytical support to national law enforcement agencies.\footnote{INTERPOL. (2023). \emph{Financial crimes strategy}. \url{https://www.interpol.int/}} There was also the creation of the Egmont Group of Financial Intelligence Units, who facilitate the exchange of financial intelligence between national agencies.\footnote{Egmont Group. (2023). \emph{Annual report 2022–2023}. \url{https://egmontgroup.org/}}

\subsection{Possible Solutions}

The solution with the broadest consensus is to strengthen transparency of beneficial ownership.\footnote{van der Does de Willebois, E., Halter, E. M., Harrison, R. A., Park, J. W., \& Sharman, J. C. (2011). \emph{The puppet masters: How the corrupt use legal structures to hide stolen assets and what to do about it}. World Bank.} This would register across all jurisdictions. The part without consensus is whether these registers should be publicly accessible, such as in the United Kingdom. Or whether it should be restricted to a closed list of law enforcement authorities.\footnote{Findley, M. G., Nielson, D. L., \& Sharman, J. C. (2014). \emph{Global shell games: Experiments in transnational relations, crime, and terrorism}. Cambridge University Press.}

Another possible solution is regulating professional enablers. This would mean extending mandatory obligations, such as increased customer due diligence, suspicious activity reporting, and record-keeping requirements. These obligations would extend to lawyers, accountants, trust and company service providers, real estate agents and dealers in high-value goods.\footnote{Halliday, T., Levi, M., \& Reuter, P. (2014). \emph{Global surveillance of dirty money: Assessing assessments of regimes to control money-laundering and combat the financing of terrorism}. American Bar Foundation.} \footnote{Boles, J. R. (2015). The two faces of bribery: International corruption pathways meet conflicting legislative regimes. \emph{Michigan Journal of International Law, 35}(4), 673–713.} This would require overcoming professional lobbies which are often very powerful, and would require much negotiation between oversight for regulations and the privilege portion of a lawyers job.\footnote{European Commission. (2021). \emph{A package of legislative proposals to strengthen the EU's anti-money laundering and countering the financing of terrorism rules}. \url{https://ec.europa.eu/}}

Investigating and researching challenges for regulating emerging financial technologies is a less straight-forward process.\footnote{Financial Action Task Force. (2021). \emph{Updated guidance for a risk-based approach to virtual assets and virtual asset service providers}. \url{https://www.fatf-gafi.org/}} This extends to cryptocurrencies and DeFi platforms and increased compliance for the use and creation of blockchain technologies. Investigations could go beyond the current technologies and move to developing new monitoring and preventative technologies, such as monitoring real-time transaction tools.\footnote{Foley, S., Karlsen, J. R., \& Putniņš, T. J. (2019). Sex, drugs, and bitcoin: How much illegal activity is financed through cryptocurrencies? \emph{Review of Financial Studies, 32}(5), 1798–1853.} \footnote{Auer, R., Haslhofer, B., Kitzler, S., Mayer, C., \& Zetzsche, D. (2023). The technology of decentralized finance (DeFi). \emph{BIS Working Papers, No. 1066}. Bank for International Settlements.}

Real-time financial intelligence sharing through INTERPOL would be an incredibly practical measure. It would require establishing secure, and standardised channels for fast exchange of financial intelligence. The main parties would be national law enforcement agencies. It could build on the Egmont group model, however, it would be advised to increase operational integration.\footnote{Egmont Group. (2023). \emph{Annual report 2022–2023}. \url{https://egmontgroup.org/}} This could dramatically decrease the time and resources required to conduct investigations.\footnote{Sheptycki, J. (2004). Organizational pathologies in police intelligence systems: Some contributions to the lexicon of intelligence-led policing. \emph{European Journal of Criminology, 1}(3), 307–332.} However, any mechanism like this would require robust safeguards, for example, against misuse, data protection standards, or independent oversight.\footnote{Deflem, M. (2010). \emph{The policing of terrorism: Organizational and global perspectives}. Routledge.}
